\section{What can go wrong, and how hard is repair?}
\begin{frame}[t,label=example]{A small example}
\small
Each client has its own firm. Every pair is feasible.
Initially $E_0=\{ay,bx,cz\}$.
\[
\begin{array}{lll}
 a:x\succ y\succ z & b:y\succ x\succ z & c:x\succ z\succ y\\
 x:b\succ c\succ a & y:a\succ b\succ c & z:c\succ a\succ b
\end{array}
\]
\vspace{-2mm}
\begin{columns}[T,onlytextwidth]
\begin{column}{.31\textwidth}\centering
No audits\par\matchpic{1}\par Stable
\end{column}
\begin{column}{.36\textwidth}\centering
Audit $a,b$\par\matchpic{2}\par {\color{rust}$cx$ blocks}
\end{column}
\begin{column}{.31\textwidth}\centering
Also audit $c$\par\matchpic{3}\par Stable again
\end{column}
\end{columns}
\sources{Solid lines are assignments. The dashed pair is hidden.
All three stability checks use the same full graph and the same preferences.}
\end{frame}

\begin{frame}[t,label=chain]{The rejection chain}
\small
After auditing $a,b$, both get their first choices. But $c$ still only sees $z$.
Now $c$ prefers $x$, and $x$ prefers $c$ to its new partner $a$.

\medskip
After auditing $c$, restart DA. Let $a,b$ propose first, then $c$:
\begin{center}
\begin{tabular}{@{}lll@{}}
\toprule
Proposal & Listing keeps & Displaced client\\
\midrule
$c$ to $x$ & $c$ & $a$\\
$a$ to $y$ & $a$ & $b$\\
$b$ to $x$ & $b$ & $c$\\
$c$ to $z$ & $c$ & none\\
\bottomrule
\end{tabular}
\end{center}
We return to the original stable matching.
Adding audits can turn success into failure. Success is not monotone.
\end{frame}

\begin{frame}[t,label=gap]{Two audits versus $N+2$}
\small
\textbf{Theorem.} For every $N\geq1$, unit costs can give
$\supp=2$ and $\resp=N+2$.

Take $N$ copies of the example. All $a_i$ share firm $F_A$, all $b_i$
share $F_B$, and other firms are private. Each copy starts at its stable matching.
\begin{center}
\begin{tikzpicture}[x=1cm,y=1cm,font=\fontsize{10}{12}\selectfont]
\node (shared) at (5.3,1.4) {Audit $F_A,F_B$: cost $2$};
\node[draw=ink,align=center,minimum width=3cm,minimum height=.85cm] (one) at (1.5,0)
 {Copy $1$: $c_1x_1$ blocks\\one private repair};
\node[draw=ink,align=center,minimum width=3cm,minimum height=.85cm] (two) at (5.1,0)
 {Copy $2$: $c_2x_2$ blocks\\one private repair};
\node at (7.25,0) {$\cdots$};
\node[draw=ink,align=center,minimum width=3cm,minimum height=.85cm] (last) at (9.4,0)
 {Copy $N$: $c_Nx_N$ blocks\\one private repair};
\draw[-{Stealth[length=1.6mm]},blue] (shared.south)--(5.3,.75)-|(one.north);
\draw[-{Stealth[length=1.6mm]},blue] (5.3,.75)-|(two.north);
\draw[-{Stealth[length=1.6mm]},blue] (5.3,.75)-|(last.north);
\end{tikzpicture}
\end{center}
Each shared firm also represents $N+3$ buyers in separate hidden one-edge
markets. Their sellers have private firms. These make substitutes for $F_A,F_B$ expensive.
\sources{Unit costs and side-exclusive firms. The shared portfolios grow with $N$.}
\end{frame}

\begin{frame}[t,label=gap-proof]{Why the gap is exact}
\small
\begin{enumerate}
\setlength{\itemsep}{6pt}
\item \textbf{Support costs two.} $F_A,F_B$ open every isolated pair.
The original stable matchings remain available. One audit cannot open
both isolated-pair families.
\item \textbf{A cheapest repair uses both shared firms.} Omitting $F_A$
requires $N+3$ private seller audits, plus one for $F_B$'s family:
at least $N+4$. The same holds for omitting $F_B$.
\item \textbf{Each copy costs one more.} With $F_A,F_B$ audited,
we must audit $f(c_i)$ or $f(x_i)$. Either exposes $c_ix_i$ and repairs that copy.
\end{enumerate}
Those $N$ private choices are disjoint. Auditing $F_A,F_B$ and all $f(c_i)$
achieves $N+2$, so the bound is tight.
\sources{Appendix 1 gives a gap of $1$ versus $N+1$ with one client per firm.
This is a cost separation, not an approximation-hardness theorem.}
\end{frame}

\begin{frame}[t,label=switch]{A two-by-two switch}
\small
For an edge $e=\{u,v\}$ in a graph, use buyers $p,q$ represented by firms
$F_u,F_v$. Sellers $x,y$ have private firms. Initially expose $px,qx$.
\[
 p:y\succ x,\quad q:x\succ y,\quad
 x:p\succ q,\quad y:q\succ p.
\]
\begin{center}
\begin{tabular}{@{}lll@{}}
\toprule
Relevant audit & DA outcome & Full-market stability\\
\midrule
None & $px$, with $q,y$ unmatched & $qy$ blocks\\
$F_u$ & $py,qx$ & stable\\
$F_v$ alone & $px,qy$ & stable\\
$y$'s firm & $py,qx$ & stable\\
\bottomrule
\end{tabular}
\end{center}
Success is exactly $F_u\in A$ or $F_v\in A$ or $f(y)\in A$.
\end{frame}

\begin{frame}[t,label=hardness]{A reduction from Vertex Cover}
\small
\setlength{\parskip}{2pt}
Vertex Cover selects vertices touching every edge. Use degree $\leq3$.
\par\smallskip
{\centering
\begin{tabular}{@{}ll@{}}
\toprule
Graph & Matching instance\\
\midrule
Vertex $v$ & Buying firm $F_v$\\
Edge $\{u,v\}$ & One switch using firms $F_u,F_v$\\
Degree of $v$ & Number of clients represented by $F_v$\\
\bottomrule
\end{tabular}
\par}
\smallskip
A cover audits an endpoint of every switch. To go back, replace private
$y$-firm audits by endpoint firms. The OR condition preserves success
without increasing cost.

The optima coincide: NP-complete budget decision, APX-hard optimization.
This holds with at most three clients per firm and two-item lists.
\sources{Unit costs and side-exclusive firms. Verification is polynomial.
Degree-three Vertex Cover: Garey et al.\ (1976), Alimonti and Kann (2000).}
\end{frame}


